
\documentclass{article}
\usepackage{colortbl}
\usepackage{makecell}
\usepackage{multirow}
\usepackage{supertabular}

\begin{document}

\newcounter{utterance}

\centering \large Interaction Transcript for game `cladder', experiment `full\_v1.5\_default', episode 8089 with score{-}Qwen3.5{-}9B{-}sp.
\vspace{24pt}

{ \footnotesize  \setcounter{utterance}{1}
\setlength{\tabcolsep}{0pt}
\begin{supertabular}{c@{$\;$}|p{.15\linewidth}@{}p{.15\linewidth}p{.15\linewidth}p{.15\linewidth}p{.15\linewidth}p{.15\linewidth}}
    \# & \multicolumn{2}{c}{Player} && \multicolumn{2}{c}{Game Master} \\
    \hline

    \theutterance \stepcounter{utterance}  
    & & & \multicolumn{4}{p{0.6\linewidth}}{
        \cellcolor[rgb]{0.9,0.9,0.9}{
            \makecell[{{p{\linewidth}}}]{
                \texttt{\tiny{[P1$\langle$GM]}}
                \texttt{You are an expert in causal inference. The following question is not a typical commonsense query, but rather a meticulously designed question created by a professor specializing in causal inference, intended to assess the students' mastery of the course content.} \\
\texttt{Imagine a self{-}contained, hypothetical world with only the following conditions, and without any unmentioned factors or causal relationships: Kwox has a direct effect on kwoz. Swoq has a direct effect on kwoz. The overall probability of kwox is 1\%. For those who are not kwox, the probability of kwoz is 47\%. For those who are kwox, the probability of kwoz is 92\%. Is kwoz less likely than not kwoz overall?} \\
\texttt{Start your answer with 'Yes' or 'No', followed by additional reasoning or evidence to support your explanation.} \\
            }
        }
    }
     \\ \\

    \theutterance \stepcounter{utterance}  
    & \multicolumn{4}{p{0.6\linewidth}}{
        \cellcolor[rgb]{0.561,0.737,0.561}{
            \makecell[{{p{\linewidth}}}]{
                \texttt{\tiny{[P1$\rangle$GM]}}
                \texttt{No} \\
\\ 
\texttt{**Reasoning:**} \\
\\ 
\texttt{To determine if \$kwoz\$ is less likely than not \$kwoz\$ overall, we must calculate the total probability of \$kwoz\$ (\$P(kwoz)\$) using the Law of Total Probability. We are given the following conditions:} \\
\\ 
\texttt{1.  **Prior Probability of Kwox:** \$P(kwox) = 0.01\$ (1\%).} \\
\texttt{2.  **Complement Probability:** Consequently, the probability of not being \$kwox\$ is \$P(\textbackslash{}neg kwox) = 1 {-} 0.01 = 0.99\$ (99\%).} \\
\texttt{3.  **Conditional Probability of Kwoz given Kwox:** \$P(kwoz | kwox) = 0.92\$ (92\%).} \\
\texttt{4.  **Conditional Probability of Kwoz given Not Kwox:** \$P(kwoz | \textbackslash{}neg kwox) = 0.47\$ (47\%).} \\
\\ 
\texttt{Now, we calculate the total probability of \$kwoz\$:} \\
\texttt{\$\$P(kwoz) = P(kwoz | kwox) \textbackslash{}cdot P(kwox) + P(kwoz | \textbackslash{}neg kwox) \textbackslash{}cdot P(\textbackslash{}neg kwox)\$\$} \\
\texttt{\$\$P(kwoz) = (0.92 \textbackslash{}cdot 0.01) + (0.} \\
            }
        }
    }
    & & \\ \\

    \theutterance \stepcounter{utterance}  
    & & & \multicolumn{2}{p{0.3\linewidth}}{
        \cellcolor[rgb]{0.9,0.9,0.9}{
            \makecell[{{p{\linewidth}}}]{
                \texttt{\tiny{[GM$|$GM]}}
                \texttt{No} \\
            }
        }
    }
    & & \\ \\

    \theutterance \stepcounter{utterance}  
    & & & \multicolumn{2}{p{0.3\linewidth}}{
        \cellcolor[rgb]{0.9,0.9,0.9}{
            \makecell[{{p{\linewidth}}}]{
                \texttt{\tiny{[GM$|$GM]}}
                \texttt{yes} \\
            }
        }
    }
    & & \\ \\

    \theutterance \stepcounter{utterance}  
    & & & \multicolumn{2}{p{0.3\linewidth}}{
        \cellcolor[rgb]{0.9,0.9,0.9}{
            \makecell[{{p{\linewidth}}}]{
                \texttt{\tiny{[GM$|$GM]}}
                \texttt{game\_result = LOSE} \\
            }
        }
    }
    & & \\ \\

\end{supertabular}
}

\end{document}
